\documentclass[11pt,a4paper]{article}
\usepackage[margin=3.2cm]{geometry}
\usepackage[T1]{fontenc}
\usepackage[utf8]{inputenc}
\usepackage{parskip}
\usepackage{enumitem}

% Long \texttt{} paths and identifiers cannot hyphenate; let TeX stretch the
% interword space instead of running them into the margin.
\sloppy
\emergencystretch=3em

\title{Retrieval + VLM Gate --- Setup and Run Guide}
\author{Nishank Satish}
\date{\today}

\begin{document}
\maketitle

This is a step-by-step guide to getting the retrieval half of the project
running from a fresh clone. Follow the steps in order --- each one depends
on the one before it. Nothing here needs you to understand embeddings or
retrieval; for what the system actually does, see \texttt{RETRIEVAL\_README.md}
in the project root instead.

\section{What you need before starting}

\begin{itemize}[nosep]
  \item This repository, on the correct branch.
  \item An NVIDIA API key (free) for the VLM step.
  \item The \texttt{data} folder --- sent separately on the shared Google
    Drive, not in git.
  \item Python 3.10 or later, and \texttt{pip}.
\end{itemize}

\section{Step 1 --- Get the code}

Clone the repository and switch to the branch with the retrieval work on
it. \texttt{main} and \texttt{feat/object-detection-YOLO} do not have any
of this --- it is all on its own branch:

\begin{verbatim}
git clone \
  git@github.com:terriljoel/retrieval-grounded-remote-sensing.git
cd retrieval-grounded-remote-sensing
git checkout feat/retrieval-vlm-gate
\end{verbatim}

\section{Step 2 --- Python environment}

Two requirements files, one for each half of the project:

\begin{verbatim}
pip install -r requirements.txt
pip install -r requirements-retrieval.txt
\end{verbatim}

\section{Step 3 --- API key}

The VLM step (the second AI model that reads uncertain detections) calls
NVIDIA's hosted API. Get a free key at
\texttt{https://build.nvidia.com} (any model page, then ``Get API Key''), then:

\begin{verbatim}
cp .env.example .env
\end{verbatim}

Open \texttt{.env} and fill in:

\begin{verbatim}
NIM_API_KEY=nvapi-...your key...
NIM_MODEL=nvidia/nemotron-3-nano-omni-30b-a3b-reasoning
\end{verbatim}

The \texttt{NIM\_MODEL} line matters --- the code's built-in default model
was retired by NVIDIA and will fail with an error if you skip this.

\texttt{.env} is never read directly by the code and is never committed to
git. Before running anything that calls the VLM, load it into your shell:

\begin{verbatim}
export NIM_API_KEY='nvapi-...'
export NIM_MODEL='nvidia/nemotron-3-nano-omni-30b-a3b-reasoning'
\end{verbatim}

or, equivalently, \texttt{set -a; . ./.env; set +a}.

\section{Step 4 --- Place the dataset}

\textbf{The NWPU VHR-10 dataset} (satellite images with labelled objects)
is the only raw dataset anything in this guide needs. It goes at:

\begin{verbatim}
docs/shared_resources/datasets/raw/NWPU VHR-10 dataset/
    positive image set/
    negative image set/
    ground truth/
\end{verbatim}

These exact folder names matter --- the code reads this layout directly.
Ask for this dataset if you were not sent it separately; nothing in the
project runs without it.

(HRRSD is a second dataset that shows up elsewhere in this repository, but
nothing in this guide needs it --- it is used by one optional script that
already ran and produced its output, and neither the demo app nor any of
the measurement scripts in Step 6 touch it. Skip it.)

\textbf{The \texttt{data} folder} you receive on the shared Drive replaces
(or merges into) the project's own \texttt{data/} directory at the
repository root:

\begin{verbatim}
data/
    memory/       <- verified-case database. Required for the demo app.
    crops/        <- cached crop images. Optional.
    embeddings/   <- cached embeddings. Optional.
\end{verbatim}

Only \texttt{data/memory/} is strictly required (for the demo app in Step
7). The others are caches --- every script rebuilds them automatically if
they are missing, just more slowly on first run.

\section{Step 5 --- Model weights}

Nothing to do here. The RemoteCLIP encoder (about 2.3~GB) downloads itself
automatically from Hugging Face the first time any script runs, as long as
you have an internet connection. It is cached afterwards under
\texttt{checkpoints/} and will not re-download.

\section{Step 6 --- Run the measurement scripts}

These reproduce every number in \texttt{reports/FINDINGS.md}. Each one is
self-contained --- run from the repository root:

\begin{verbatim}
python3 -m scripts.run_gate_test     # does retrieval work at all?
python3 -m scripts.run_baselines     # the two simple baselines
python3 -m scripts.run_alpha_sweep   # tuning the crop fusion weight
python3 -m scripts.run_open_set      # spotting an unseen class
python3 -m scripts.run_label_noise   # robustness to bad labels
python3 -m scripts.run_arm_ablation  # which crop views help
python3 -m scripts.run_conformal     # calibrating the accept threshold
python3 -m scripts.run_b3_vlm \
    --scoring logprob --evidence none
\end{verbatim}

Each writes its results to \texttt{reports/gate/} or \texttt{reports/phase1/}
and prints a summary to the terminal. The first run of each is slower
(building crops and embeddings); later runs read from cache and are quick.

\section{Step 7 --- The real detector export (optional, but needed for the demo)}

Two things in this project need a real YOLO detection export, not just the
dataset: the interactive demo app, and \texttt{scripts/run\_b0.py} (the
real-vs-synthetic comparison). Place the export at:

\begin{verbatim}
data/detections/detections.csv
data/detections/inference_images.csv   <- if your export has one
\end{verbatim}

If you do not have this file, ask for it, or generate your own by running
the detector's own export step (see the main \texttt{README.md} for that
half of the project).

\section{Step 8 --- Run the demo app}

\begin{verbatim}
streamlit run app/gate_explorer.py
\end{verbatim}

Open the URL it prints. Pick an image, pick a detection, click ``Find
similar cases,'' click ``Ask the VLM,'' and see the gate's verdict next to
the ground truth. \texttt{docs/UI\_TEST\_CASES.md} has four specific cases
worth clicking through, with what each one is supposed to show.

\section{If something goes wrong}

\begin{itemize}[nosep]
  \item \textbf{A script says a file is not found under \texttt{docs/shared\_resources}} ---
    the dataset is not in the right place; recheck Step 4's folder names exactly.
  \item \textbf{The VLM step raises an error naming \texttt{NIM\_API\_KEY}} ---
    the key is not loaded into your shell; see Step 3.
  \item \textbf{The VLM step raises an HTTP error about a model} ---
    \texttt{NIM\_MODEL} is not set, or is set to the retired default; see Step 3.
  \item \textbf{The app says \texttt{detections.csv} is missing} --- see Step 7.
  \item \textbf{Anything else} --- the error messages in this codebase are
    written to say exactly what is missing and where it was expected; read
    the message before asking, it usually says the fix directly.
\end{itemize}

\section{Where to read more}

\begin{description}
  \item[\texttt{RETRIEVAL\_README.md}] What this half of the project does,
    the eight-step pipeline explained in plain English, and the headline
    results.
  \item[\texttt{reports/FINDINGS.md}] Every result from this project,
    numbered, tagged by whether it is measured on real data or synthetic
    proposals only.
  \item[\texttt{reports/gate/GATE\_REPORT.md}] The first test of whether
    retrieval works at all.
  \item[\texttt{docs/UI\_TEST\_CASES.md}] Four worked examples for the demo
    app, with what each one is supposed to show.
  \item[\texttt{docs/STATE.md}] The full running log of decisions and
    findings, if you want the complete history rather than the summary.
\end{description}

\end{document}
